% TEMPLATE EXEMPLAR of a crash-course chapter: it uses every element once (opening, goals, sections, run-in
% paragraphs, equation, Key idea, Example, figure, table, Caution, Connection, Build it with a tested listing,
% Check yourself, Further reading). Copy its shape, then replace it with the book's first chapter (git mv it to the
% new name and overwrite it). Its labels are also used in main.tex, chapters/preface.tex,
% chapters/appendix_glossary.tex and figures/fig_chapterdeps.tex; find them all with
%   grep -rn -E 'ch:(tour|reading)|part:(start|literature)|fig_decay' main.tex chapters figures
% Labels: <type>:<chapter slug>-<name>, e.g. sec:tour-rc, fig:tour-decay (tools/check.sh enforces the types).
\chapter{Exponential Decay in One Sitting}
\label{ch:tour}

A capacitor emptying through a resistor, the charge of a battery left on a shelf, the weight that a running
average gives to old samples, the step size of many learning-rate schedules: all of them shrink the same way. Each
loses the same \emph{fraction} of what is left in every equal stretch of time, and that one rule gives a curve you
will meet in every chapter that follows. This chapter takes one sitting. It needs nothing beyond the exponential
function from high-school algebra. By the end you should be able to

\begin{itemize}
  \item say what a time constant and a half-life are, and convert one into the other;
  \item predict how much of a decaying quantity is left after any time, without a calculator for the common cases;
  \item compute the time constant of a resistor and a capacitor, and of a running average.
\end{itemize}

\section{One rule, one curve}
\label{sec:tour-curve}

A quantity \emph{decays exponentially} when, in every interval of the same length, it loses the same fraction of
its current value. If a quantity $v$ starts at $V_0$, the rule gives
\begin{equation}
  v(t) = V_0\, e^{-t/\tau},
  \label{eq:tour-decay}
\end{equation}
where $\tau$ (the Greek letter tau) is the \emph{time constant}, measured in seconds. After one time constant
$e^{-1}\approx 0.37$ of the start is left, so the quantity has lost \qty{63}{\percent}.

\paragraph{Half-life.} The \emph{half-life} $t_{1/2}$ is the time after which half is left. Setting
$e^{-t/\tau}=1/2$ gives $t_{1/2}=\tau\ln 2\approx0.69\,\tau$: the half-life is a little shorter than the time
constant, and both describe the same curve (\cref{fig:tour-decay}).

\begin{insightbox}[The decay rule]
\[
  v(t)=V_0\,e^{-t/\tau},\qquad t_{1/2}=\tau\ln 2\approx 0.69\,\tau .
\]
After $\tau$, $2\tau$, $3\tau$ and $5\tau$, about \qty{37}{\percent}, \qty{14}{\percent}, \qty{5}{\percent} and
\qty{0.7}{\percent} of the start is left (\cref{tab:tour-fractions}).
\end{insightbox}

% Example of a computed figure: the curve is plotted from the formula in the text (eq:tour-decay) with the
% numbers of the worked example (V0 = 5 V, tau = RC = 1 s), never drawn by hand.
\begin{figure}[tbp]
\centering
\begin{tikzpicture}
\begin{axis}[bookplot,width=0.74\linewidth,height=4.8cm,xmin=0,xmax=5.2,ymin=0,ymax=5.6,
  xlabel={time $t$ (s)},ylabel={voltage $v$ (V)},
  xtick={0,1,2,3,4,5},ytick={0,1,2,3,4,5},clip=false]
  \addplot[navy,domain=0:5,samples=120] {5*exp(-x)};
  % one time constant: 37 percent left
  \addplot[only marks,mark=*,mark size=2pt,rose] coordinates {(1,1.8394)};
  \draw[rose,dashed,line width=0.5pt] (axis cs:1,0) -- (axis cs:1,1.8394) -- (axis cs:0,1.8394);
  \node[font=\sffamily\scriptsize,text=rose,anchor=west] at (axis cs:1.12,2.05) {$t=\tau$: \qty{1.84}{V}, \qty{37}{\percent} left};
  % half-life: 50 percent left
  \addplot[only marks,mark=*,mark size=2pt,teal] coordinates {(0.6931,2.5)};
  \node[font=\sffamily\scriptsize,text=teal,anchor=west] at (axis cs:0.80,3.05) {$t_{1/2}=\tau\ln 2=\qty{0.69}{s}$: half left};
  % three time constants: 5 percent left
  \addplot[only marks,mark=*,mark size=2pt,amber] coordinates {(3,0.2489)};
  \node[font=\sffamily\scriptsize,text=amber,anchor=south west] at (axis cs:3.05,0.35) {$t=3\tau$: \qty{0.25}{V}, \qty{5}{\percent} left};
\end{axis}
\end{tikzpicture}
\caption{A \qty{100}{\micro\farad} capacitor charged to \qty{5}{V} discharging through \qty{10}{\kilo\ohm}
($\tau=RC=\qty{1}{s}$), computed from \cref{eq:tour-decay}. Every interval of $\tau$ removes the same fraction,
\qty{63}{\percent} of what is left. Original drawing.}
\label{fig:tour-decay}
\end{figure}


\section{A capacitor that empties}
\label{sec:tour-rc}

A \emph{capacitor} stores electric charge; the voltage across it is proportional to the charge it holds. Connect a
charged capacitor $C$ (in farads, \unit{\farad}) across a resistor $R$ (in ohms, \unit{\ohm}) and the current that
flows through the resistor drains the charge. The more charge is left, the higher the voltage and the faster the
current, so the voltage falls by the same fraction in every equal interval: it decays exponentially, with time
constant $\tau=RC$ \cite{horowitz2015art}.

\begin{workedbox}[How long until it is empty?]
A \qty{100}{\micro\farad} capacitor charged to \qty{5}{V} discharges through \qty{10}{\kilo\ohm}. The time
constant is
\[
  \tau = RC = \qty{10000}{\ohm}\times\qty{100e-6}{\farad} = \qty{1.0}{s}.
\]
After \qty{1}{s} the voltage is $5\times e^{-1}=5\times0.368=\qty{1.84}{V}$; after \qty{3}{s}, $5\times e^{-3}=\qty{0.25}{V}$. To
fall to \qty{1}{V}, solve $5\,e^{-t/\tau}=1$: $t=\tau\ln 5=\qty{1.6}{s}$. A decaying quantity never reaches zero,
so ``empty'' always means ``below a threshold you choose''; five time constants (\qty{0.7}{\percent} left) is a
common choice.
\end{workedbox}

\begin{table}[tbp]
\centering
\small
\caption{The fraction $e^{-t/\tau}$ left after whole multiples of the time constant, and the voltage of the
worked example ($V_0=\qty{5}{V}$, $\tau=\qty{1}{s}$).}
\label{tab:tour-fractions}
\begin{tabular}{@{}rrr@{}}
\toprule
$t$ & fraction left & voltage \\
\midrule
$\tau$  & \qty{36.8}{\percent} & \qty{1.84}{V} \\
$2\tau$ & \qty{13.5}{\percent} & \qty{0.68}{V} \\
$3\tau$ & \qty{5.0}{\percent}  & \qty{0.25}{V} \\
$5\tau$ & \qty{0.67}{\percent} & \qty{0.03}{V} \\
\bottomrule
\end{tabular}
\end{table}

\begin{cautionbox}[A charged capacitor can bite]
A capacitor keeps its charge after the circuit is switched off, and large ones, such as those inside power supplies,
can hold enough energy to burn or shock. Do not open a power supply. On your own low-voltage circuits, discharge a
large capacitor through a resistor of a few hundred ohms and check with a meter before you touch it; never short it
with a screwdriver.
\end{cautionbox}

\begin{connectionbox}[Running averages decay too]
An exponential moving average $\bar x_t=\beta\,\bar x_{t-1}+(1-\beta)\,x_t$ weights a sample $k$ steps old by
$(1-\beta)\beta^k$. For $\beta$ close to 1, $\beta^k\approx e^{-k(1-\beta)}$, so the weights decay with a time
constant of about $1/(1-\beta)$ steps: $\beta=0.99$ averages over roughly the last 100 steps.
\end{connectionbox}

\begin{buildbox}[Compute the curve yourself]
This Python snippet (\texttt{figures/src/tour\_decay.py}; its output is saved next to it, and
\texttt{tools/check.sh} reruns it) tabulates \cref{eq:tour-decay} for the worked example.
\lstinputlisting[language=Python]{figures/src/tour_decay.py}
It prints \texttt{5.00}, \texttt{3.03}, \texttt{1.84}, \texttt{0.68} and \texttt{0.25}~V, and a half-life of
\texttt{0.693}~s.
\end{buildbox}

\begin{checkbox}
(1)~A quantity has $\tau=\qty{2}{s}$. What fraction is left after \qty{6}{s}? (2)~A half-life is \qty{10}{min}.
What is $\tau$? (3)~What is the time constant of a running average with $\beta=0.9$?

{\footnotesize Answers: (1)~$6/2=3$ time constants, $e^{-3}\approx\qty{5}{\percent}$. (2)~$\tau=t_{1/2}/\ln2
\approx\qty{14.4}{min}$. (3)~About $1/(1-0.9)=10$ steps.\par}
\end{checkbox}

\begin{readingbox}[Further reading]
The first chapter of Horowitz and Hill~\cite{horowitz2015art} treats resistor-capacitor circuits and their time
constant, with the engineer's rules of thumb. This chapter left out the mirror image, exponential approach to a
final value ($v=V_f+(V_0-V_f)e^{-t/\tau}$), which follows from the same rule.
\end{readingbox}
